\documentclass{article}
\usepackage{amsmath,amssymb}
\usepackage[margin=1in]{geometry}
\allowdisplaybreaks
\usepackage[UTF8]{ctex}
\begin{document}

\section*{非参数化的个体-客户MDP近似方法}

对于每一个当前状态$x\in\mathcal X$，请选择相应的策略。
\begin{equation}
\hat\pi(x)\in\operatorname*{arg\,max}_{a\in\mathcal A(x)}
\left\{g(x,a)+\alpha\int_{\mathcal X}\hat V(y)\,Q[dy\mid x,a]\right\}.
\tag{6}
\end{equation}
续接的值是非线性背包问题的最优解。
\begin{equation}
\begin{aligned}
\hat V(x)&\equiv\max_{k_n}\sum_{n=1}^{N}V_n^*(x_n,k_n)\\
\text{s.t.}\quad&\sum_{n=1}^{N}t_nk_n\leq MT.
\end{aligned}
\tag{5}
\end{equation}
这里$k_n$表示可供客户$n$使用的车辆行程数量，因此是一个非负整数。打印块(5)中没有为$k_n$单独设置领域行或上限值。固定数量$t_n$表示从供应商到客户$n$再返回所需的时间；$M$表示车辆数量；$T$表示每辆车每个时间段的最大工作时间。背包算法构建的是一种价值，而不是(6)中描述的可行行动集：论文明确指出，只有当这种分配对于IRPDD来说是可行的，它才被视为一种可能的政策。

\section*{继承的输入、状态和可行的路线}

有一个产品，一个供应商的工厂，以及由$n=1,\ldots,N$索引的$N$个客户。车队拥有$M$辆具有已知容量$C_V$的同类车辆。客户$n$的存储能力为$C_n$。时间是以离散单位计算的，$t=0,1,\ldots$，而一个周期被称为一天。库存状态为：
\[
\mathcal X=[0,C_1]\times[0,C_2]\times\cdots\times[0,C_N],\qquad
X_{nt}\in[0,C_n],\qquad
X_t=(X_{1t},\ldots,X_{Nt})\in\mathcal X.
\]
在第6.1节中的分解部分，论文假设需求分布以及状态空间$\mathcal X$是离散的。不过，这里并没有明确说明具体的离散化网格结构。

供应商已知客户需求的固定联合分布$F$，且该分布不会随时间发生变化。不同日期上的需求向量是相互独立的；同一天内不同客户之间的需求也不具有独立性。用$U_{nt}$表示客户$n$在日期$t$的需求，用$F_n$表示$U_n$的边际分布，用$f_n(u_n)$表示客户$n$的需求为$u_n$的概率（以离散形式表示）。折扣因子为：
\[
\alpha\in[0,1).
\]
供应商可以在任何时间测量客户的库存水平。该行程的持续时间少于一天，因此每天开始时，所有的车辆和司机都处于可用状态。一辆车每天可以执行多条路线。

对于直接配送专享服务，$\mathcal A(x)$指的是那些每条路线只服务一个客户的行程安排。这些行程安排需要满足继承下来的工作量、行驶时间、司机工作时间限制、客户交付时间窗口、车辆容量、客户交付后的存储能力以及其他特定于应用的行程限制。这些限制是在行程安排中明确规定的，而不是作为单独的代数式来表述的。

一个行程$a$由各个客户的行程$a_n$和$n=1,\ldots,N$组成。每个$a_n$都规定了每辆车对客户$n$的访问次数，以及每辆车所交付的数量。数量$d_n(a)$，或者直接用直接交付的表示法表示为$d_n(a_n)$，是指交付给客户$n$的总量；$k_n(a_n)$则表示车辆访问该客户的次数。函数$k_n(a_n)$用于计算实际的访问次数，而状态/背包变量$k_n$则表示可提供的行程次数。

继承下来的可行性要求包括：
\[
A_t\in\mathcal A(X_t)\quad(t=0,1,\ldots),
\qquad
X_{nt}+d_n(A_t)\leq C_n
\quad(n=1,\ldots,N,\ t=0,1,\ldots).
\]
后者是作者实际使用的版本，假设在库存测量和交付之间没有使用任何产品。任何路线上都不能超过车辆容量。方程(5)中的总工作量不平衡并不替代这些路线限制。

每个动作$A_t$必须仅依赖于历史信息。
\[
(X_0,A_0,X_1,\ldots,X_t),
\]
不要使用未来的信息。一个静态的确定性策略$\pi\in\Pi$仅根据当前状态$x$来选择$\pi(x)\in\mathcal A(x)$。其客户$n$的行程则由$\pi_n(x)$表示。

\section*{库存动态与系统奖励}

未满足的需求将被视为已丢失，不会积压。对于每个客户和每一天，所使用的数量、短缺量以及下一批库存的数量都将被记录下来。
\[
\min\{X_{nt}+d_n(A_t),U_{nt}\},
\]
\[
S_{nt}=\max\{U_{nt}-(X_{nt}+d_n(A_t)),0\},
\qquad
X_{n,t+1}=\max\{X_{nt}+d_n(A_t)-U_{nt},0\}.
\]
对于$x\in\mathcal X$、$a\in\mathcal A(x)$以及可测量的$B\subseteq\mathcal X$，该论文进行了定义。
\[
\mathcal U(x,a,B)\equiv
\left\{U\in\mathbb R_+^N:
\max\{x_n+d_n(a)-U_n,0\}\in B\right\},
\]
\[
P[X_{t+1}\in B\mid X_t=x,A_t=a]
=Q[B\mid x,a]\equiv F[\mathcal U(x,a,B)].
\]
$\mathcal U$的定义显示保留在其印刷形式的符号表示中，尽管$B$是向量状态空间的一个子集。

固定的函数$r_n(d_n)$、$h_n(\cdot)$和$p_n(s_n)$分别对应于交付的奖励、库存持有成本以及未满足需求的惩罚。它们仍然只是抽象的函数而已。固定的$c_n$则是向客户$n$进行交付所需的运输成本。在(6)中使用的直接交付预期的单阶段系统奖励就是……
\begin{equation}
\begin{aligned}
g(x,a)\equiv\sum_{n=1}^{N}\Bigl\{
&r_n(d_n(a_n))-c_nk_n(a_n)-h_n(x_n+d_n(a_n))\\
&-E^{F_n}\!\left[p_n\!\left(\max\{U_n-(x_n+d_n(a_n)),0\}\right)\right]
\Bigr\}.
\end{aligned}
\tag{3}
\end{equation}
这里的$E^{F_n}$表示相对于边际需求分布$F_n$的期望值。所有数字都应保持不变。

\section*{单个客户之间的转换构造}

固定策略$\pi\in\Pi$，并假设在该策略下存在唯一的静态概率$\nu^\pi(x)$。在当前的库存状况$x_n$和交付情况$d_n$下，下一个库存到达客户$n$的概率$y_n$，以及下一个策略行程会访问该客户$k_n$次的概率，就是这样的。
\begin{equation}
q_n(y_n,k_n\mid x_n,d_n)=
\frac{
\displaystyle\sum_{\{s\in\mathcal X:\ s_n=x_n,\ d_n(\pi_n(s))=d_n\}}
\nu^\pi(s)
\displaystyle\sum_{\{z\in\mathcal X:\ z_n=y_n,\ k_n(\pi_n(z))=k_n\}}
Q[z\mid s,\pi(s)]
}{
\displaystyle\sum_{\{s\in\mathcal X:\ s_n=x_n,\ d_n(\pi_n(s))=d_n\}}
\nu^\pi(s)
}.
\tag{4}
\end{equation}
此表达式适用于分母为正的情况。打印出的特殊情况为：
\[
q_n(y_n,k_n\mid x_n,d_n)=0
\quad\text{if the denominator in (4) is zero}.
\]
对于这种情况，没有提供任何替代的概率分布。

对于每一个$n=1,\ldots,N$，客户MDP的状态为$(x_n,k_n)$：$x_n$是客户的库存，而$k_n$是可供使用的车辆次数。库存坐标$x_n,y_n$属于$[0,C_n]$范围内的相应离散库存状态空间；$k_n,l_n$则表示访问次数。可接受的方案$\mathcal A(x_n,k_n)$包括向客户$n$派遣最多$k_n$次车辆，而交付量则受到车辆容量$C_V$和客户存储容量$C_n$的限制。

其打印出的转换关系为：
\[
\begin{aligned}
&P[(X_{n,t+1},K_{n,t+1})=(y_n,l_n)
\mid (X_{nt},K_{nt})=(x_n,k_n),A_{nt}=a_n]\\
&\hspace{35mm}=q_n(y_n,l_n\mid x_n,d_n(a_n)).
\end{aligned}
\]
第6.1节中规定的单阶段奖励是
\[
\begin{aligned}
g_n(x_n,a_n)\equiv{}&r_n(d_n(a_n))-c_nk_n(a_n)-h_n(x_n)\\
&-E^{F_n}\!\left[p_n\!\left(\max\{U_n-(x_n+d_n(a_n)),0\}\right)\right].
\end{aligned}
\]
其最优的折现价值是
\[
V_n^*(x_n,k_n)\equiv
\sup_{\{A_{nt}\}_{t=0}^{\infty}}
E\!\left[
\left.\sum_{t=0}^{\infty}\alpha^t g_n(X_{nt},A_{nt})
\right|(X_{n0},K_{n0})=(x_n,k_n)
\right],
\]
受…的约束
\[
A_{nt}\in\mathcal A(X_{nt},K_{nt}),
\qquad n=1,\ldots,N,\quad t=0,1,\ldots,
\]
并且，这些操作必须是非预期的：它们仅使用在决策时可用的客户流程历史数据。这些就是(5)中使用的个体客户值。

\paragraph{打印奖励矛盾。}
第6.1节，物理页面11上指出：
\[
g(x,a)=\sum_{n=1}^{N}g_n(x_n,a_n).
\]
不过，$g_n$的定义中包含了$-h_n(x_n)$，而(3)中则包含了$-h_n(x_n+d_n(a_n))$。因此，这种等价关系并不符合这两个印刷定义的一致性。在上述情况下，由于(6)中的系统奖励采用了(3)中的定义，而第6.1节中的印刷定义则被用于客户价值的计算。因此，这种等价关系被记录为源矛盾，而不是被当作额外的身份或用来改变任何假设条件。

\section*{选定的非参数估计与策略更新}

所选版本为算法1，使用的是物理页面14中简化的条件概率。这里采用了以下方法来代替直接的五维经验更新（7）；两个估计量并不是同时施加的。

在固定的模拟策略$\pi$下，定义客户$n$被车辆$k_n$访问的概率，这一概率取决于库存状况$y_n$。
\[
p_n(k_n\mid y_n)=
\frac{
\displaystyle\sum_{\{x\in\mathcal X:\ x_n=y_n,\ k_n(\pi_n(x))=k_n\}}\nu^\pi(x)
}{
\displaystyle\sum_{\{x\in\mathcal X:\ x_n=y_n\}}\nu^\pi(x)
}
\quad\text{if the denominator is positive}.
\]
以不同的客户索引来打印零分母句子：
\[
p_i(k_i\mid y_i)=0\quad\text{if the denominator is zero}.
\]
这部分内容保留了印刷版的特殊情况表示法。条件概率$p_n(k_n\mid y_n)$与短缺惩罚函数$p_n(s_n)$是不同的；论文中使用了相同的字母来表示两者，只是通过它们的参数来区分。

让$\hat p_{nt}(k_n\mid y_n)$在模拟经过$t$次转换后估计出$p_n(k_n\mid y_n)$的值；同时，让$N_{nt}(y_n)$表示客户$n$在通过转换$t$后处于库存状态$y_n$的次数。初始估计值用$N_{n0}(k_n\mid y_n)$表示，这与式(7)中的初始观测值类似。打印出来的更新结果如下：
\[
\hat p_{n,t+1}(k_n\mid y_n)=
\begin{cases}
\displaystyle
\frac{(N_{n0}(k_n\mid y_n)+N_{nt}(y_n))\hat p_{nt}(k_n\mid y_n)+1}
{N_{n0}(k_n\mid y_n)+N_{nt}(y_n)+1},
&\begin{array}{l}
\text{if }X_{nt}=y_n\text{ and}\\
k_n(\pi_n(X_t))=k_n,
\end{array}\\[4mm]
\displaystyle
\frac{(N_{n0}(k_n\mid y_n)+N_{nt}(y_n))\hat p_{nt}(k_n\mid y_n)}
{N_{n0}(k_n\mid y_n)+N_{nt}(y_n)+1},
&\begin{array}{l}
\text{if }X_{nt}=y_n\text{ and}\\
k_n(\pi_n(X_t))\ne k_n,
\end{array}\\[4mm]
\hat p_{nt}(k_n\mid y_n),&\text{if }X_{nt}\ne y_n.
\end{cases}
\]
客户转换概率的估计结果为：
\[
\hat q_{nt}(y_n,k_n\mid x_n,d_n)=
\begin{cases}
f_n(x_n+d_n-y_n)\hat p_{nt}(k_n\mid y_n),&\text{if }y_n>0,\\[1mm]
\displaystyle\sum_{u_n=x_n+d_n}^{\infty}f_n(u_n)\hat p_{nt}(k_n\mid y_n),&\text{if }y_n=0.
\end{cases}
\]
下方的端点$u_n=x_n+d_n$以及单独的边界情况$y_n=0$保持不变。论文指出，这些估计值通常与使用公式(7)得到的数值有所不同。方程(4)定义了被近似的条件概率；并不是说，经过简化的估计值必须等于那个静态概率比率。在更新的客户MDP中，估计的转移概率取代了上述转移关系中的确切数值$q_n$。

论文指出，$\hat p$的估计结果具有与直接经验性估计相似的收敛性：对于频繁出现的状态，在固定的模拟策略下，这些估计结果以概率一收敛。该构造中并未规定具体的数值初始权重、估计值或停止容忍度。

\paragraph{算法1：策略更新构建方法。}
\begin{enumerate}
\item 以初始策略$\hat\pi_0$为起点，然后设置$i=0$。
\item 重复以下模拟、解决方案、策略更新和增量操作，进行选定的迭代次数，或直到满足收敛测试条件为止。
\item 在 $\hat\pi_i$ 条件下模拟 IRPDD 的运行，以使用上述选定的更新来估算 $p_n(k_n\mid y_n)$ 的值。
\item 根据更新后的估算结果，重新构建并解决更新后的个体客户MDP，从而得到$V_n^*(x_n,k_n)$。
\item 通过(6)定义$\hat\pi_{i+1}$，并使用那些更新的客户值，通过(5)定义$\hat V$。
\item 增加$i\leftarrow i+1$的值。
\end{enumerate}

\paragraph{来源范围。}
本模型采用 Kleywegt、Nori 和 Savelsbergh 的《\emph{直接配送下的随机库存路径问题}》（1999 年 6 月）第 2、5、6.1 和 6.2 节，见 PDF 实际第 4 至 6 页和第 10 至 15 页。 参数近似、基于采样的动作比较程序、贪婪动作子问题以及附录A中的CBW公式均不属于本选段内容。 没有提供额外的官方补充材料；任何额外的补充资料均未经过验证。

\end{document}
